\section*{\S~2b Wahl der Leitung} 
\label{sec:wahl_derleitung}
\addcontentsline{toc}{section}{\S~2b Wahl der Leitung}

\begin{enumerate}[label=(\arabic*)]
    \item Die Leitung wird durch die Mitgliederversammlung mit einer Zweidrittelmehrheit für die Dauer von einem Kalenderjahr gewählt und befördert.
    \item Es gibt keine Obergrenze für Amtszeiten. Auch können beliebig viele Amtszeiten hintereinander erfolgen.
    \item Um für einen Leitungsposten kandidieren zu können, muss ein Mitglied mindestens 24 Monate Mitglied gewesen sein und darf am Wahltag nicht länger als vor 12 Monate zuletzt ein passives Mitglied gewesen sein. Zudem muss das Mitglied die Gruppenführer-Kampfrolle, eine abgeschlossene Zeus/Missionsbauer-Ausbildung und eine weitere Sonder-Kampfrolle besitzen. Eine vorherige Mitarbeit im Stab ist empfehlenswert.
    \item Sollte der Unit-Leiter vor Ende seiner Amtszeit nicht mehr in der Lage sein, sein Amt auszuführen, übernimmt der Spieß die Leitung der Unit. Der Posten des Unit-Leiters bleibt bis zur nächsten Wahl unbesetzt. Diese sollte zum nächst möglichen Termin durchgeführt werden.
    \item Sollte der Spieß vor Ende seiner Amtszeit nicht mehr in der Lage sein, sein Amt auszuführen, bleibt der Posten des Spieß bis zur nächsten Wahl unbesetzt. Diese sollte zum nächst möglichen Termin durchgeführt werden.
    \item Die Amtszeit beginnt am 01. Januar eines Jahres und endet am 31. Dezember des gleichen Jahres. Wird im Verlauf eines Kalenderjahres eine andere Leitung gewählt, so ist dennoch am Termin beschrieben in \S~2e Absatz 4 \ref{itm:regularer_Termin} eine Wahl durchzuführen.
\end{enumerate}
